% Pass5316--5321: Latin/Q4 derivative and D4/tomotope square-face refinement.
\subsection{Affine Latin colorings, $D_4$ triality, and the tomotope square-face action}

Fix the explicit toroidal-knight--to--$Q_4$ labeling used by the router
certificate, and encode the four Latin symbols by $\mathbb F_2^2$.  Pulling an
order-four Latin square back to $\mathbb F_2^4$, define the four first
differences
\[
  \Delta_iL(x)=L(x+e_i)+L(x)\in\mathbb F_2^2.
\]
An exhaustive census of all $576$ labeled Latin squares gives a sharp affine
criterion.  The $144$ Klein-$V_4$ squares are exactly those for which all four
$\Delta_iL$ are constant, equivalently the symbol map
$L:\mathbb F_2^4\to\mathbb F_2^2$ is affine.  None of the $432$ cyclic-$C_4$
squares is affine: $192$ have exactly one constant derivative ($48$ for each
of the four $Q_4$ directions) and $240$ have none.  This refines the fixed
knight-edge census of Pass5311; it is a coloring/code statement, not a symmetry
group identification.

For the Hoffman $13$-cell cover, let $H$ be the order-$576$ stabilizer and let
$N\cong W(D_4)$ be the canonical order-$192$ normal subgroup of Pass5308, so
$H/N\cong C_3$.  The ten $H$-orbits on four-cell subsets from Pass5302 refine
to exactly $22$ $N$-orbits.  Four $H$-orbits remain transitive under $N$ and
six split into triples of equal $N$-orbits.  Thus the outer $C_3$ acts literally
as a triality fusion on the first open four-cell shortening layer.  This orbit
refinement does not determine the shortened minimum distance.

Under the Pass5300 affine isomorphism
\[
  H/Z(H)\cong L_+,
\]
where $L_+$ is the order-$288$ even-parastrophe group of the Klein Latin
square, the characteristic order-$96$ subgroup corresponding to $N/Z(N)$
fixes the row, column, and symbol lines of the five-line $PG(3,2)$ spread
individually and swaps the two orthogonal-mate lines.  The quotient $C_3$ is
represented by cyclic parastrophe,
\[
  \mathrm{row}\longmapsto\mathrm{column}\longmapsto\mathrm{symbol}
  \longmapsto\mathrm{row},
\]
while fixing both mate lines.  This is the Latin-square realization of the
outer $D_4$ triality factor; it is distinct from the inner $S_3$ quotient of
the tomotope group.

There is also a natural tesseract realization of the tomotope action.  The
$tesseract$ has $24$ square faces, hence $12$ antipodal square-face pairs.
The even-sign signed-permutation group $W(D_4)$ acts on those $12$ objects,
and its central antipode is exactly the kernel.  The induced order-$96$ action
is permutation-conjugate to the published Monson--Pellicer--Williams
tomotope action; an explicit conjugating relabeling is
\[
  (0,1,4,5,8,9,11,10,7,6,3,2).
\]
Consequently
\[
  W(D_4)/\langle -I\rangle\curvearrowright
  \{12\ \text{antipodal square-face pairs}\}
  \cong \Gamma(\mathcal T)
\]
as degree-$12$ permutation representations.

The three small non-diagonal orbitals of this action are perfect matchings.
Their union is $3K_4$.  The three $K_4$ components are exactly the three
complementary $2+2$ partitions of the four tesseract coordinate axes,
\[
  01\mid23,\qquad 02\mid13,\qquad 03\mid12.
\]
The kernel of the induced action on these three components is elementary
abelian of order $16$, while the component action is $S_3$.  Hence the natural
square-face action realizes the semidirect product directly:
\[
  \Gamma(\mathcal T)\cong (C_2)^4:S_3.
\]

Finally, the orientation-preserving tesseract group $R$ and $W(D_4)$ remain
distinct order-$192$ double covers.  On the same $12$ antipodal square-face
pairs, $W(D_4)/Z$ has orbital rank $5$ (three separate perfect-matching
relations), whereas $R/Z$ has orbital rank $4$: it preserves one matching and
fuses the other two into a degree-two relation.  Both share the same uncolored
$3K_4$ skeleton, but the natural degree-$12$ representations are not
permutation-conjugate.  No physical identification is inferred from these
finite symmetry correspondences.
